% ===== §13 Return and the Spiral =====
\section{Return and the Spiral}
\label{sec:return}

\noindent
A system driven far from equilibrium fluctuates, and the previous section established that estrangement is the signature of that fluctuation. But fluctuation is not mere departure. A dissipative system that persists does not simply wander away from its states; it \emph{returns} to them---and the entire normative weight of the paper rests on a distinction in the \emph{form} of that return. There is a return that comes back to exactly where it began, and a return that comes back lifted; and the difference between them is the difference between a relation that is dying and one that is alive. This section makes that distinction precise, drawing only on the dynamical and dialectical resources already in play. It is the hinge between the thermodynamics of the driven system and the dormancy that the next section treats.

\subsection{Two forms of return}
\label{sub:return-two}

Consider a system that cycles. After a phase of departure---of estrangement, of withdrawal, of the disequilibrium the last section described---it comes back toward coupling. The coming-back can take one of two forms, and they are distinguished exactly by the holonomic criterion the series has used throughout (\xref{sec:phase}).

The first is the \emph{closed circle}: a return to the identical configuration from which the cycle set out, having traversed a loop that accumulated no net phase. This is the return of zero holonomy. The system passes through its phases and arrives precisely where it started, neither lifted nor depleted; nothing has been generated by the circuit. A relation whose every cycle of estrangement-and-reconciliation returns it to exactly the prior state---the same quarrel, the same repair, the same equilibrium, endlessly---is executing closed circles. It is not dead, for it still moves; but it generates nothing by its motion, and a system that cycles without accumulating is, as the cosmological frame insists, on its way to the equilibrium it merely postpones.

The second is the \emph{spiral}: a return that comes back to a configuration \emph{like} its starting point but displaced---lifted along an axis the planar circle does not have, carrying a surplus that the cycle generated and that was not present at the outset. This is the return of positive holonomy. The system passes through the same \emph{kind} of phase it passed through before---another estrangement, another withdrawal, another return---but arrives not at the identical state but at a raised one, enriched by the traverse. The estrangement was not wasted; the loop, in closing, lifted.

\begin{claim}[The spiral, not the circle]
The return of a generative relation is a \emph{spiral}, not a circle. The closed circle returns to the identical state, accumulating no phase (zero holonomy): the same cycle of rupture and repair endlessly repeated, generating nothing by its motion. The spiral returns to a state \emph{like} the prior one but lifted, carrying the surplus the cycle generated (positive holonomy): the same kind of estrangement traversed, but issuing in a raised configuration the planar circle has no axis to reach. The criterion distinguishing a living return from a dying one is therefore not whether the relation cycles---every driven system cycles---but whether its cycling is helical or planar: whether each return lifts, or merely repeats.
\end{claim}

\subsection{The spiral is the negation of the negation}
\label{sub:return-dialectic}

This is, in the dialectical vocabulary of the preceding section, simply the \emph{negation of the negation}. The first negation is the departure---the estrangement, the contradiction, the phase of disequilibrium that negates the prior coupling. The second negation does not restore the prior coupling unchanged, as though the departure had not happened; it negates the negation, and in doing so produces a new coupling that contains the departure as a moment sublated within it. The result is not the original state recovered but a higher state attained \emph{through} the rupture---which is precisely the spiral, and precisely the \emph{Aufhebung} the account of the union of two material histories reached by its own road (\xref{sub:histmat-sublation}). The dialectical spiral and the positive-holonomy cycle are not two doctrines but one structure described in two vocabularies: the return-with-surplus that distinguishes a generative cycle from a sterile one.

It follows that estrangement is, for a spiralling relation, not the enemy of return but its \emph{means}. The lift is generated \emph{through} the departure, not despite it; a relation that suffered no estrangement would have no first negation to negate, and so no spiral---only the flat persistence of an unchallenged state, which is the closed circle in its limiting, motionless case. This is the strongest form the paper's first claim can take. Estrangement is not merely tolerable, not merely a necessary cost; it is the raw material of the lift, the negation without which there is no negation of the negation, the departure without which there is no enriched return.

\subsection{Why the spiral cannot be commanded}
\label{sub:return-cannot}

One caution preserves the consistency of this section with the methodological core (\xref{sec:conditions}). That the return \emph{can} be a spiral does not mean the spiral can be commanded. Whether a given departure issues in an enriched return or in mere repetition is not a matter the parties can directly will, for it depends on whether the loop closes generatively---an event of the real, not a deployment of the symbolic. What the parties can do is tend the conditions under which a spiral is possible rather than foreclosed: keep the channel clear enough that the departure's errors remain attributable (\xref{sec:prediction}), refrain from depositing the negative curvature that makes the next loop harder to lift (\xref{sub:phase-forms}), guard against the freezing that converts a living spiral into a closed circle (\xref{sec:balance}). The spiral is not produced by these; it is made possible by them. One cannot command the lift, only keep open the conditions under which lifting may occur---which is, once more, the whole discipline of the political economy of conditions, here in its temporal aspect.

\begin{claim}[The spiral is conditioned, not commanded]
That a return can be a spiral does not mean the spiral can be willed. Whether a departure issues in an enriched return or in sterile repetition depends on the generative closing of the loop, an event of the real beyond direct symbolic command. The parties can only tend the conditions under which a spiral is possible---a channel clear enough for error to remain attributable, a geometry not freshly scarred, a power not frozen into a circle. The lift is conditioned, not commanded; the most one can do is keep open the conditions under which the return may lift rather than merely repeat.
\end{claim}

\noindent
The spiral, then, names the form of a living return: a coming-back enriched through the very departure that seemed to threaten the bond. But a spiral has phases of descent as well as ascent, and there are stretches of its arc in which the system appears not to be generating at all---in which it has withdrawn, gone quiet, drawn inward. Read against the closed circle, such a phase looks like failure; read against the spiral, it may be something else entirely: the low, inward phase from which the next ascent is gathered. To that phase---the dormancy that is not decline but the storing of force for return---the final cosmological section turns.
